\section{\ref{sec:sensors}장. 기계의 입력}

센서의 구조는 단일 세포 전류 기록, 달팽이관 진동 기록, 망막 세포 계수로 직접 측정되었다. 감도의 한계와 산탄 잡음 식은 물리에서 나온다. 센서의 부호가 최대 정보에 맞춰져 있다는 것은 강한 근거를 가진 가설이다.

{\footnotesize
\setlength{\tabcolsep}{3pt}\renewcommand{\arraystretch}{1.15}
\begin{longtable}{>{\raggedright}p{0.5cm}>{\raggedright}p{4.0cm}>{\raggedright}p{5.6cm}>{\raggedright}p{2.3cm}>{\raggedright\arraybackslash}p{1.7cm}}
\toprule
No. & 명제 & 실험과 측정한 것 & 출처 & 상태 \\
\midrule\endfirsthead
\toprule
No. & 명제 & 실험과 측정한 것 & 출처 & 상태 \\
\midrule\endhead
1 & 센서는 변환기이다: 수용체가 전류를 만들고, 눈과 귀 감각 세포의 출력은 \nt{GLUT} 버스로 가는 글루탐산이다; 첫 세포들은 스파이크를 내지 않는다(그림~\ref{fig:transducer}) & 거북의 막대세포와 원뿔세포는 흥분성 아미노산을 방출한다: 떼어 낸 막 조각의 글루탐산 채널이 이를 잡는다. 쥐 달팽이관의 속털세포: 신경섬유 말단의 전류는 글루탐산 AMPA 수용체를 거치며, 방출 빈도는 세포의 완만한 탈분극과 함께 $150$\,Hz까지 커진다 & \cite{CopenhagenJahr1989,GlowatzkiFuchs2002} & 측정됨 \\
2 & 어둠 속 빛 센서는 광자 하나에 약 1피코암페어의 전류로 응답한다 & 두꺼비 막대세포 바깥 분절에 흡입 전극: 희미한 섬광에 대한 응답이 양자화되어 있고, 사건은 $\approx1$\,pA(암전류의 $5\%$), 퍼짐 $0.2$\,pA, 효율 $0.5$. 사람은 광자 하나가 들어온 것을 우연보다 자주 알아맞힌다 & \cite{Baylor1979,Tinsley2016} & 측정됨 \\
3 & 소리 센서는 1나노미터보다 작은 변위를 알아챈다 & 뫼스바우어 방법, 기니피그 달팽이관 기저막의 $16$--$19$\,kHz 위치: 청신경 응답 문턱에서 막의 속도는 약 $0.04$\,mm/s, 곧 변위 $\approx0.35$\,nm(우리의 환산). 잰 것은 센서 다발이 아니라 막이다 & \cite{Sellick1982,Hudspeth1989} & 측정됨 \\
4 & 어둠 속 정밀도는 광자 산탄 잡음~\eqref{eq:photonnoise}이 제한한다; 약한 빛에서는 적분이 길어진다 & 식은 푸아송 통계에서 나온다. 막대세포에서 희미한 빛의 잡음은 독립적인 양자 사건들의 합과 일치한다; 밝은 배경에서 섬광 응답은 더 작고 짧다. “수십분의 1초”라는 수치는 여기서 따로 검증되지 않음 & \cite{Baylor1979} & 이론 \\
5 & 입력 범위는 빛이 열 자릿수, 소리가 열두 자릿수이고, 스파이크 발화율은 세 자릿수가 안 된다 & 소리: 특성 주파수에서 기저막 응답은 $40$--$80$\,dB 대역에서 기울기 $0.2$\,dB/dB까지로 커지며, 압축은 $100$\,dB 위에서도 보인다. 자릿수 자체는 표준 추정치로, 본문에서는 출처 없이 제시 & \cite{Ruggero1997} & 측정됨 \\
6 & 센서 응답은~\eqref{eq:nakarushton}이며, $\sigma$는 평균 밝기를 따라가 곡선을 로그 축을 따라 이동시킨다(그림~\ref{fig:adaptation}) & 이 식은 처음 물고기 망막의 S 전위에 맞춰졌다. 거북 원뿔세포: 응답은 $V=I V_m/(I+\sigma)$를 따르고 밝기 $\sim3.5$자릿수를 담는다; 배경에 $2$--$3$분 노출하면 곡선은 폭을 유지한 채 수평으로 이동한다. 전체 활동에 의한 나눗셈과의 관계는 총설 & \cite{NakaRushton1966,NormannPerlman1979,Carandini2012} & 측정됨 \\
7 & 센서는 변화를 전하고, 그 부호는 스파이크당 최대 정보에 맞춰져 있다(명제~\ref{prop:sensors}) & 원리는 이론적으로 제안되었다. 가장 강한 증거: 파리에서 첫 층 뉴런의 “대비 $\to$ 응답” 특성이 자연 장면 대비의 누적 분포와 일치한다. 이렇게 해서 최대 정보가 얻어진다 & \cite{Barlow1961,Laughlin1981} & 가설 \\
8 & 켜짐 센서와 꺼짐 센서는 두 전선 부호이다; 이벤트 카메라는 이를 실리콘으로 재현한다 & 실험 총설: 망막의 켜짐 채널과 꺼짐 채널은 첫 시냅스에서 이미 갈라지고 따로 차단된다. 카메라: 프레임 없는 $128\times128$픽셀, 범위 $120$\,dB, 지연 $15\,\mu\text{s}$ & \cite{Schiller1992,Lichtsteiner2008,Gallego2022} & 측정됨 \\
9 & 눈에는 빛 센서 $\sim10^8$개와 출력 뉴런 $\sim10^6$개가 있다: $\sim100$배 압축 & 사람 망막 전체 표본에서의 계수: 평균 막대세포 $92$백만 개, 원뿔세포 $4.6$백만 개; 출력(신경절) 세포는 눈에 따라 $0.7$에서 $1.5$백만 개 & \cite{Curcio1990,CurcioAllen1990} & 측정됨 \\
10 & 생쥐 눈의 출력 채널은 서른 개가 넘는다 & 생쥐: 출력 세포 $11\,000$개 이상의 이광자 칼슘 영상과 응답의 군집화, $30$개보다 뚜렷이 많은 기능 유형. 사람의 수는 측정되지 않음 & \cite{Baden2016} & 측정됨 \\
11 & 기니피그의 눈은 $\sim875\,000$ bit/s를 전송한다; 사람 눈으로 환산하면 $\sim10^7$ bit/s & 기니피그, 다전극 배열, 자연 장면 속 움직임: 세포당 $6$--$13$ bit/s, 출력 세포 $\sim10^5$개에 $\sim875\,000$ bit/s. 사람($\sim10^6$개 세포)의 $\sim10^7$ bit/s는 저자들의 환산 & \cite{Koch2006} & 측정됨 \\
12 & 귀는 거의 로그인 주파수 지도를 가진 대역 통과 필터 뱅크이다(그림~\ref{fig:filterbank}) & 기저막 최대 진동 위치는 주파수에 따라 다르다; “주파수--위치” 함수는 거의 지수적이며, 하나의 형태로 사람과 살아 있는 동물의 달팽이관을 기술한다 & \cite{Greenwood1990,Robles2001} & 측정됨 \\
13 & 공진기는 능동적이다: 센서 일부가 약한 진동을 진폭으로 수천 배($66$--$76$\,dB) 증폭하고, 소리가 커지면 증폭이 줄어든다 & 친칠라 달팽이관 기저부의 레이저 진동 측정: 약한 순음의 특성 주파수에서 등자뼈 대비 이득 $66$--$76$\,dB; 죽으면 감도가 $60$--$81$\,dB(진폭으로 $10^3$--$10^4$배) 떨어지고 압축이 사라진다. 힘의 원천은 바깥털세포 & \cite{Ruggero1997,Robles2001} & 측정됨 \\
14 & 귀의 증폭기는 자기 발진 경계 근처에 있다 & 계산: 호프 분기점 근처에서는 범위 압축, 날카로운 조율, 결합음이 불가피하다. 이것이 알려진 청각의 비선형성 전부이다. 달팽이관이 실제로 그렇게 조율되어 있다는 직접 증명은 없음 & \cite{Eguiluz2000} & 가설 \\
15 & 낮은 주파수에서 스파이크는 소리와 위상을 맞춰 나가며(기니피그에서 위상 고정은 $\sim600$\,Hz부터 약해지고 $\sim3.5$\,kHz까지 드러난다), 이로부터 방향을 찾는다; 높은 주파수에서는 위치 부호만 남는다 & 기니피그 청신경: 위상 고정은 약 $600$\,Hz에서 떨어지기 시작해 $3.5$\,kHz 위에서 사라진다; 고양이와 원숭이에서는 경계가 약 한 옥타브 높다. 두 귀의 도착 시간 차는 총설 & \cite{PalmerRussell1986,Grothe2010} & 측정됨 \\
16 & 나머지 센서(표~\ref{tab:sensors}): 후각은 수용체 유형 $\sim400$개, 뉴런당 하나, 조합 부호; 미각은 일부 ATP와 세로토닌을 통해; 촉각은 빨리, 느리게 적응하는 센서; 따뜻함과 차가움은 별도 & 생쥐: 한 수용체가 많은 냄새를, 한 냄새가 많은 수용체를 인식하며, 냄새마다 집합이 다르다. ATP 수용체가 없으면 미각 신경이 응답하지 않는다; 미뢰는 세로토닌을 방출한다. 사람 손 피부의 단일 섬유 기록: 적응 속도에 따른 네 유형. 냉각 채널은 열 채널과 다르다 & \cite{Mombaerts2004,Malnic1999,Finger2005,Huang2005,JohanssonVallbo1979,McKemy2002} & 측정됨 \\
\bottomrule
\end{longtable}}
